\documentclass[11pt]{article}
\usepackage[margin=1in]{geometry}
\usepackage{amsmath,amssymb,amsfonts,mathtools}
\usepackage{array,booktabs,enumitem,graphicx,xcolor,hyperref}
\usepackage[T1]{fontenc}
\usepackage[utf8]{inputenc}
\title{Novel Symmetries of Topological Conformal Field theories}
\author{J. Sonnenschein and S. Yankielowicz}
\date{}
\begin{document}
\maketitle
\begin{abstract}

We show that various actions of topological conformal theories that were suggested recentely are particular cases of a general action. We prove the invariance of these models under transformations generated by nilpotent fermionic generators of arbitrary conformal dimension, \$\textbackslash{}Q\$ and \$\textbackslash{}G\$. The later are shown to be the \$n\textasciicircum{}\{th\}\$ covariant derivative with respect to ``flat abelian gauge field" of the fermionic fields of those models. We derive the bosonic counterparts \$\textbackslash{}W\$ and \$\textbackslash{}R\$ which together with \$\textbackslash{}Q\$ and \$\textbackslash{}G\$ form a special \$N=2\$ super \$W\_\textbackslash{}infty\$ algebra. The algebraic structure is discussed and it is shown that it generalizes the so called ``topological algebra".

\end{abstract}

\section{Introduction}

W\textasciicircum{} (n) Q\textasciicircum{} (n) R\textasciicircum{} (n) G\textasciicircum{} (n) \_ m+n,0 2cm \#1 Comm. Math. Phys. \#1 \#1 Class. Quantum Grav. \#1 \#1 Phys. Lett. \#1B \#1 Phys. Rev. Lett. \#1 \#1 Phys. Rev. D\#1 \#1 Phys. Rev. \#1 \#1 Nucl. Phys. B\#1 \#1 Nuovo Cimento \#1 \#1 J. Math. Phys. \#1 \#1 Mod. Phys. Lett. A\#1 E. Witten, 117 (1988) 353. E. Witten, 340 (1990) 281. R. Dijkgraaf, E. Witten, 342 (1990) 486.

\section{Method}

R. Dijkgraaf, E. Verlinde, and H. Verlinde, 352 (1991) 59; ``Notes On Topological String THeory And 2D Quantum Gravity,'' Princeton preprint PUPT-1217 (1990). K. Aoki, D. Montano, J. Sonnenschein, 247 (1990) 64. D. Montano, J. Sonnenschein, 313 (1989) 258. J. Labastida, M. Pernici, E. Witten, 310 (1988) 611. D. Montano, J. Sonnenschein , 324 (1989) 348. J. Sonnenschein 42 (1990) 2080. T. Eguchi and S.-K. Yang, 5 (1990) 1693. K. Li, 354 (1991) 711, ibid. 725

\begin{equation}
\label{eq:compact}
L(\theta) = \sum_{i=1}^{n} \ell(x_i, \theta)
\end{equation}
Equation~\ref{eq:compact} is included to keep the sample paper-like while remaining small.

\section{Conclusion}

E. Bergshoeff, C.N. Pope, L. J. Romans, E. Sezgin, X. Shen, and S. Stelle 243 (1990) 350, 245 (1990) 447, 256 (1990) 350; `` math expression Gravity and Super math expression CERN-TH 5743/90; E. Bergshoeff, M. Vasiliev and B. de Wit 256 (1991) 199 and references therein. K. Aoki, D. Montano, J. Sonnenschein ``The role of Contact Algebra in the Multi Matrix Models" Caltech preprint CALT-68-1676, to appear in Mod. Phys. Lett. D. Friedan, E. Martinec and S. Shenkar 271 (1986) 93
\end{document}
